\section{Computing and Auditing Fisher-Glass}
\label{sec:computing-auditing}

An operational audit fixes the task contrast, nuisance hypothesis, matrix construction, and audit distribution before measuring information on independent data. We work in low-dimensional task and nuisance subspaces, so the reported information concerns these declared coordinates. Coordinate normalization is specified for each experiment; fitted nuisance probes and temperature calibration are frozen before final evaluation. The audit reports support and conditioning alongside its information values.

\subsection{Operational Scores and Nuisance Coordinates}
\label{sec:operational-scores}

The audit uses a declared matrix target. Predictive Fisher is computed from likelihood scores under the fitted predictive distribution; observed-label score moments and unweighted feature moments define additional audit targets. The prospective Waterbirds experiment uses the uncentered feature--Schur fraction of fixed task-logit and background-probe coordinates. Hidden-slice discovery uses estimated temperature-calibrated predictive Fisher with class-specific shrinkage $\eta=0.10$, retained during search and held-out evaluation. Unweighted feature geometry supplies a complementary comparison. Appendix~\ref{app:operational-matrices} specifies each construction.

The experiments prespecify nuisance coordinates from metadata, fitted annotation probes, or a fixed semantic direction. Without nuisance annotations, the specified semantic direction serves as the audited nuisance hypothesis. We report fitting splits, normalization, and nuisance-block conditioning; cross-validation assesses gate stability. Appendix~\ref{app:predefined-diagnostics} records the checks performed. Formal anchor task-preservation tests and admission thresholds are criteria for extending this family, rather than implemented selection steps.

\subsection{Stable Estimation and Exact-Information Audits}
\label{sec:stable-estimation}

Let $M_i$ be the per-example task--nuisance matrix under the chosen construction, with task and nuisance blocks indexed by $T$ and $N$. For a predefined slice $g$ of size $n_g$, form $\widehat J_g=n_g^{-1}\sum_{i\in g}M_i$ and compute the ridge Schur estimate
\begin{equation}
\widehat F_g^{\rid}
=\widehat J_{TT,g}-\widehat J_{TN,g}
(\widehat J_{NN,g}+\lambda_{\rid}I)^{-1}\widehat J_{NT,g}.
\label{eq:ridge-estimator}
\end{equation}
Real-data computations use $\lambda_{\rid}=10^{-3}$. Ridge stabilizes residualization but raises information relative to the unregularized target. Let $a_g$ bound estimation error and $d_g$ bound the ridge-to-unregularized gap on a simultaneous event over the audited slices. For declared weights $\Pi_G$, use
\begin{equation}
\begin{aligned}
\underline\lambda_g&=[\lambda_{\min}(\widehat F_g^{\rid})-a_g-d_g]_+,\\
\underline I_{\alpha,G}&=
\left[\CVaR_{\alpha,g\sim\Pi_G}
\left((\underline\lambda_g+\varepsilon)^{-1}\right)\right]^{-1}.
\end{aligned}
\label{eq:lower-certificate}
\end{equation}
Here $[x]_+=\max\{x,0\}$. On that event, this lower-bounds the unregularized population certificate for the declared matrix and audit distribution. Certification also requires a positive lower confidence bound on nuisance conditioning. The implemented scalar audits use Bonferroni-adjusted Gaussian/Wald moment simulation at nominal 95\% confidence, conditional on the fitted model and evaluated slices. Appendix~\ref{app:estimation-certification} records the gap identity, error allowances, and numerical comparison.

\subsection{Hidden-Slice Discovery and Held-Out Certification}
\label{sec:hidden-slice-method}

A gate $h_\phi(x,y)\in[0,1]$ assigns membership weights. Discovery maximizes inverse projected information subject to effective-sample-size, maximum-weight, and class-mass constraints, with shrinkage toward the same-class pooled matrix. Five validation folds assess transfer and stability; after fixing the protocol, a final gate is fitted on all discovery data, frozen, and evaluated on an untouched split. Its unshrunk target is the weighted ratio $J_\phi=\E[h_\phi M]/\E[h_\phi]$, with positive membership mass. Uncertainty must account for numerator, normalization, and any constraint-induced weight adjustment; the reported Waterbirds target retains the stated shrinkage. Appendix~\ref{app:hidden-slice-discovery} specifies the constraints, objective, and held-out target.

\begin{algorithm}[t]
\caption{Fisher-Glass audit}
\label{alg:fisher-glass-audit}
\begin{algorithmic}[1]
\State \textbf{Specify:} predictor, matrix, coordinates, normalization, calibration, and audit weights.
\State \textbf{Discover and freeze:} predefined slices or constrained gates fitted on discovery data.
\State \textbf{Audit independently:} check support and conditioning; compute ridge, unregularized, and lower values.
\State \textbf{Report:} raw and projected information, surviving fraction, slice support, tail certificate, and uncertainty.
\end{algorithmic}
\end{algorithm}
